\documentclass{zhvt-classic}

% 使用本地源流明体（GenRyuMin2 TW）按文件路径加载，不依赖系统字体名
% 版框（内框／界栏／书口侧栏／鱼尾）改用与版心印章同一印泥棕色：
%   印章图片 yshs.png 实色 = RGB(135,68,52) = HTML 874434
\usepackage{xcolor}
\usepackage{tikz}
% sealbrown 双注册：xcolor 供 TikZ/pgf，L3 供 cls 版框
\definecolor{sealbrown}{HTML}{874434}
\ExplSyntaxOn
\color_set:nnn { sealbrown } { HTML } { 874434 }
\ExplSyntaxOff

\zhvtset{
  grid color       = sealbrown,
  main font        = GenRyuMin2TW-SB,
  main font path   = C:/Users/Administrator/texmf/fonts/opentype/public/GenRyuMin2TW/,
jiazhu font      = GenRyuMin2TW-M,
  jiazhu font path = C:/Users/Administrator/texmf/fonts/opentype/public/GenRyuMin2TW/,
  % 页码框容纳 4 字：正常页「九十七」+前置空 1 格=4 行恰好；
  % 100+ 页「一百零一」共 5 行超出，末字「一/二/三」被框裁掉（2026-10-04 实测），故放宽到 5
  shukou page number chars = 5,
  % 书口书名框高从默认 6 字减到 4 字，把省出的 2 字高拨给 chapter 框，
  % 使 6 字书名+卷名放进 9 字高的 chapter 框后不再压线（2026-10-05）
  shukou title chars = 4,
}

\title{史记}[韩长孺列传]
\author{汉·司马迁}
\maker{闲叟漫记·重製}


% ===== 索隐述赞样式（2026-11-08 四次定稿 · 愚大人亲手改定参数）=====
% ⚠ 竖排前提：自定义字体家族必须带 RawFeature={vertical:+vert}；
%   simhei.ttf 加 +vhal 会整体加载失败静默回退，只写 +vert；simfang 同理。
\newCJKfontfamily{\shuzanfont}[Path=C:/Windows/Fonts/,RawFeature={vertical:+vert}]{simhei.ttf}
\newCJKfontfamily{\shuzanbody}[Path=C:/Windows/Fonts/,RawFeature={vertical:+vert}]{simfang.ttf}
% 题头：152x30pt 圆角 6pt 棕底白字；baseline (0,2) 使整框落入界行格子不压版框线；
% 白字节点 x 不补偿、y +8.2pt（像素级墨迹实测 dx/dy≈0，L/R 留白约 20pt）。
\newcommand{\shuzanhead}[1]{%
  \begin{tikzpicture}[baseline={(0,2)}]
    \fill[sealbrown, rounded corners=6pt]
      (-76pt,-15pt) rectangle (76pt,15pt);
    \node[text=white, inner sep=0pt, anchor=center] at (0pt,8.2pt)
      {\shuzanfont\fontsize{23}{23}\selectfont #1};
  \end{tikzpicture}%
}
% ★ 修正书口卷名位置：chapter 框起始位置改回 cls 默认值，避免全局影响。
%   只在每个卷首页通过 \juan 命令里的空命令微调该页位置（安全，不会跨页污染）。
\ExplSyntaxOn
\AtBeginDocument{
  \dim_set:Nn \l__zhvt_shukou_chapter_title_start_dim
    { \l__zhvt_shukou_title_height_dim + \l__zhvt_font_size_dim * 2 }
  \dim_set:Nn \l__zhvt_shukou_page_number_start_dim
    { \l__zhvt_shukou_chapter_title_start_dim + \l__zhvt_shukou_chapter_title_height_dim }
}
\ExplSyntaxOff

\begin{document}

% ★★ 版心（书口）排布覆写：把书名放到鱼尾之下（2026-10-04）
%
%   cls 默认版心自上而下：书名（title 静态框，在鱼尾之上）→ 鱼尾 → 篇名（chapter 动态框）→ 页码
%   本模板改为：          鱼尾（位置不动）→ 书名 → 篇名 → 页码分割线 → 页码
%
%   手法：把 cls 放在 title 静态框里的书名，改放进 chapter 动态框（鱼尾之下）；
%         title 静态框留空。书名取 \l__zhvt_title_str，即 \title{...} 的第一参数，
%         所以换书时只改 \title，不必在此重复写书名。
%
%   ⚠ 三条必须遵守的写法（2026-10-04 实测）：
%     ① 汉字必须用 \hbox{...} 包住。flowfram 会把框内容装进一个宽 14.4mm（= 一个 line_sep
%        = 书口栏宽）的纵向 minipage 并按该宽度折行，而整页靠 PDF /Rotate 90 显示、
%        框内文字本身不变向 —— 不包 \hbox 时「围炉夜话」4 字会被逐字折成 4 行，
%        在版面上就成了横排 4 字并横向溢出版框。
%     ② 竖排下**一个 \hbox 就是一整列**。书名与篇名要同处一列，必须放进同一个 \hbox；
%        分成两个 \hbox 时，后一个会被当成新的一列、甩到隔壁栏去
%        （实测「序」跑到书口栏外侧 14.4mm 的正文列上）。中间的 \quad 在框内是横向空白，
%        显示出来正好是书名与篇名之间的纵向间隔。
%     ③ ★ 内容盒要写成 **与框等宽的定宽盒** `\hbox to \l__zhvt_shukou_chapter_title_height_dim{…\hfil}`。
%        实测：chapter 框是把内容盒沿纵向**居中**摆放的，内容宽度一变，书名就跟着漂：
%        自然宽度的 \hbox（书名 38.4mm）落在 x=157.2mm，加篇名后（57.6mm）落在 x=152.4mm，
%        目录页（无篇名）与正文页的书名就差半格。做成等宽盒后，无论有没有篇名，
%        内容固定从框顶起排，实测两页均落在 x=142.8mm，不再漂。
\ExplSyntaxOn
% 把书口 chapter 框（书名+篇名）起始位置下移 1 字高，避免 4 字卷名顶到鱼尾
% ⚠ cls 的 \__zhvt_calculate_dimensions: 在 \AtBeginDocument 里执行，会重算该 dim；
%   必须把覆写挂到 begin document 之后，否则会被二次覆盖（2026-10-05 实测）。
% (旧方案已移到导言区，直接修改 chapter 框起始位置)
% 卷标题命令 #1=书口篇名(鱼尾) #2=正文卷名（含卷序，卷名内空格用 ~）
%
% ★★ 目录页两项修订（2026-10-05，愚大人指出）★★
%   ① 页码起点统一对齐：
%      原状 —— 目录项写成 \par\hspace{4em}卷名\hspace{4em}页码，页码起点 = 卷名实宽 + 4em。
%      卷名实宽随字数变：「卷一~梁惠王上」6 字 +~、「卷七~离娄上」5 字 +~，
%      差 1 个字位 → 卷七~卷十 4 条的页码整体比其余 10 条高 27.21pt（一个字），看着上下不齐。
%      改法 —— 卷名补全角空格（宽 = 1 字 = 27.215pt）补足到 **7 字符位**再写进目录，
%      14 条页码起点遂全部落在同一条竖线上（起点对齐；末字随页码字数自然参差）。
%      ⚠ 只对目录里的副本 \l_tmpc_tl 补，正文卷名仍用 \l_tmpa_tl 原样 —— 正文像素不受影响。
%   ② 卷名整体上移半字位：前导 \hspace 由 4em 改 3.5em。
%      实测原状目录块在版框内偏下（天头留白 139.6pt / 地脚留白 110.7pt），
%      上移 0.5em = 13.6pt 后上下基本均分，视觉居中。
\NewDocumentCommand { \juan } { mm } {
	\clearpage
	\tl_set:Nn \l__zhvt_chapter_title_tl { #1 }
	\tl_set:Nn \l_tmpa_tl { #2 }
	\tl_set_eq:NN \l_tmpc_tl \l_tmpa_tl
	\int_step_inline:nn { 7 - \str_count:n { #2 } }
	  { \tl_put_right:Nn \l_tmpc_tl { 　 } }
	\tl_set:Nx \l_tmpb_tl { \zhnum { page } }
	\tl_put_left:Nn \l_tmpb_tl { \hspace { 4em } }
	\tl_put_left:Nn \l_tmpb_tl { \l_tmpc_tl }
	\tl_put_left:Nn \l_tmpb_tl { \par \hspace { 3.5em } }
	\addtocontents { toc } { \l_tmpb_tl }
	\skip_horizontal:n { \l__zhvt_font_size_dim * 4 }
	\l_tmpa_tl
	\par
	\thispagestyle { empty }
}
\cs_set_protected:Npn \__zhvt_set_shukou_content: {
	% ★ 鱼尾上方留白（title 静态框）嵌入闲叟钤印：
	%   整页 /Rotate 90，故此处把图用 \rotatebox{90} 转正为竖章；
	%   尺寸卡在版心条幅内（书内约 7mm 宽 × 33.8mm 高，留余量不压鱼尾）。
	\setstaticcontents * { title } {
		\vbox_to_ht:nn { \l__zhvt_baseline_skip_dim } {
			\vskip 14.2pt
			\hbox_to_wd:nn { \l__zhvt_shukou_title_height_dim } {
				\hfil
				\rotatebox [ origin = c ] { 90 } { \includegraphics [ width = 7mm , height = 33.8mm ] { yshs.png } }
				\hfil
			}
		}
	}
% ★ 版心（书口）统一改法（2026-11-16 愚大人定）：18pt＋框内居中＋书口文字同版框色
%   ① 原 27.2pt 时 7 字长篇名（如「廉颇蔺相如列传」）末字溢出 chapter 框、压住页码；
%      改 \fontsize{18pt}{18pt}，前后 \hfil 令字串在框内居中（不随篇名字数漂移）。
%   ② \vbox{\vskip 8.2pt} 定量下移，抵消缩字号后栏宽方向 8.2pt 上偏（像素实测 y 中心 543.3 = 原版）。
%      ⚠ \raisebox / \vspace* 在 flowfram 框内无效；换字号须同步换补偿：16pt→10.1、19pt→7.4。
%   ③ chapter 与 page~number 均置 \color{sealbrown}（HTML 874434），与版框／界行同色；正文仍黑。
%   ⚠ 页码字位、钤印、鱼尾、界行、正文一律未动；五部合订本同此口径。
	\setdynamiccontents * { chapter }
	{ \color{sealbrown} \vbox { \vskip 8.2pt \hbox to \l__zhvt_shukou_chapter_title_height_dim
		{ \hfil \fontsize{18pt}{18pt}\selectfont \l__zhvt_title_str \quad \l__zhvt_chapter_title_tl \hfil } } }
	\setdynamiccontents * { page~number } { \color{sealbrown} \quad \zhnum { page } }
}
\__zhvt_set_shukou_content:
\ExplSyntaxOff

\maketitle

% 单篇，无多卷目录，不渲染目录页

\juan{韩长孺列传}{卷一百零八~韩长孺列传第四十八}
\hspace*{2em}御史大夫韩安国者。梁成安人也。后徙睢阳。尝受韩子。杂家说于驺田生所。事梁孝王为中大夫。吴楚反时。孝王使安国及张羽为将。扞吴兵于东界。张羽力战。安国持重。以故吴不能过梁。吴楚已破。安国。张羽名由此显。\\
\hspace*{2em}梁孝王。景帝母弟。窦太后爱之。令得自请置相。二千石。出入游戏。僭于天子。天子闻之。心弗善也。太后知帝不善。乃怒梁使者。弗见。案责王所为。韩安国为梁使。见大长公主而泣曰。﹁何梁王为人子之孝。为人臣之忠。太后曾弗省也。夫前日吴。楚。齐。赵七国反时。自关以东皆合从西乡。惟梁最亲为艰难。梁王念太后。帝在中。而诸侯扰乱。一言泣数行下。跪送臣等六人。将兵击卻吴楚。吴楚以故兵不敢西。而卒破亡。梁王之力也。今太后以小节苛礼责望梁王。梁王父兄皆帝王。所见者大。故出称跸。入言警。车旗皆帝所赐也。即欲以侘鄙县。驱驰国中。以夸诸侯。令天下尽知太后。帝爱之也。今梁使来。辄案责之。梁王恐。日夜涕泣思慕。不知所为。何梁王之为子孝。为臣忠。而太后弗恤也。﹂大长公主具以告太后。太后喜曰。﹁为言之帝。﹂言之。帝心乃解。而免冠谢太后曰。﹁兄弟不能相教。乃为太后遗忧。﹂悉见梁使。厚赐之。其后梁王益亲驩。太后。长公主更赐安国可直千馀金。名由此显。结于汉。\\
\hspace*{2em}其后安国坐法抵罪。蒙狱吏田甲辱安国。安国曰。﹁死灰独不复然乎。﹂田甲曰。﹁然即溺之。﹂居无何。梁内史缺。汉使使者拜安国为梁内史。起徒中为二千石。田甲亡走。安国曰。﹁甲不就官。我灭而宗。﹂甲因肉袒谢。安国笑曰。﹁可溺矣。公等足与治乎。﹂卒善遇之。\\
\hspace*{2em}梁内史之缺也。孝王新得齐人公孙诡。说之。欲请以为内史。窦太后闻。乃诏王以安国为内史。\\
\hspace*{2em}公孙诡。羊胜说孝王求为帝太子及益地事。恐汉大臣不听。乃阴使人刺汉用事谋臣。及杀故吴相袁盎。景帝遂闻诡。胜等计画。乃遣使捕诡。胜。必得。汉使十辈至梁。相以下举国大索。月馀不得。内史安国闻诡。胜匿孝王所。安国入见王而泣曰。﹁主辱臣死。大王无良臣。故事纷纷至此。今诡。胜不得。请辞赐死。﹂王曰。﹁何至此。﹂安国泣数行下。曰。﹁大王自度于皇帝。孰与太上皇之与高皇帝及皇帝之与临江王亲。﹂孝王曰。﹁弗如也。﹂安国曰。﹁夫太上。临江亲父子之间。然而高帝曰﹃提三尺剑取天下者朕也﹄。故太上皇终不得制事。居于栎阳。临江王。適长太子也。以一言过。废王临江。用宫垣事。卒自杀中尉府。何者。治天下终不以私乱公。语曰。﹃虽有亲父。安知其不为虎。虽有亲兄。安知其不为狼。﹄今大王列在诸侯。悦一邪臣浮说。犯上禁。桡明法。天子以太后故。不忍致法于王。太后日夜涕泣。幸大王自改。而大王终不觉寤。有如太后宫车即晏驾。大王尚谁攀乎。﹂语未卒。孝王泣数行下。谢安国曰。﹁吾今出诡。胜。﹂诡。胜自杀。汉使还报。梁事皆得释。安国之力也。于是景帝。太后益重安国。孝王卒。共王即位。安国坐法失官。居家。\\
\hspace*{2em}建元中。武安侯田蚡为汉太尉。亲贵用事。安国以五百金物遗蚡。蚡言安国太后。天子亦素闻其贤。即召以为北地都尉。迁为大司农。闽越。东越相攻。安国及大行王恢将。未至越。越杀其王降。汉兵亦罢。建元六年。武安侯为丞相。安国为御史大夫。\\
\hspace*{2em}匈奴来请和亲。天子下议。大行王恢。燕人也。数为边吏。习知胡事。议曰。﹁汉与匈奴和亲。率不过数岁即复倍约。不如勿许。兴兵击之。﹂安国曰。﹁千里而战。兵不获利。今匈奴负戎马之足。怀禽兽之心。迁徙鸟举。难得而制也。得其地不足以为广。有其众不足以为彊。自上古不属为人。汉数千里争利。则人马罢。虏以全制其敝。且强弩之极。矢不能穿鲁缟。冲风之末。力不能漂鸿毛。非初不劲。末力衰也。击之不便。不如和亲。﹂群臣议者多附安国。于是上许和亲。\\
\hspace*{2em}其明年。则元光元年。雁门马邑豪聂翁壹因大行王恢言上曰。﹁匈奴初和亲。亲信边。可诱以利。﹂阴使聂翁壹为间。亡入匈奴。谓单于曰。﹁吾能斩马邑令丞吏。以城降。财物可尽得。﹂单于爱信之。以为然。许聂翁壹。聂翁壹乃还。诈斩死罪囚。县其头马邑城。示单于使者为信。曰。﹁马邑长吏已死。可急来。﹂于是单于穿塞将十馀万骑。入武州塞。\\
\hspace*{2em}当是时。汉伏兵车骑材官二十馀万。匿马邑旁谷中。卫尉李广为骁骑将军。太仆公孙贺为轻车将军。大行王恢为将屯将军。太中大夫李息为材官将军。御史大夫韩安国为护军将军。诸将皆属护军。约单于入马邑而汉兵纵发。王恢。李息。李广别从代主击其辎重。于是单于入汉长城武州塞。未至马邑百馀里。行掠卤。徒见畜牧于野。不见一人。单于怪之。攻烽燧。得武州尉史。欲刺问尉史。尉史曰。﹁汉兵数十万伏马邑下。﹂单于顾谓左右曰。﹁几为汉所卖。﹂乃引兵还。出塞。曰。﹁吾得尉史。乃天也。﹂命尉史为﹁天王﹂。塞下传言单于已引去。汉兵追至塞。度弗及。即罢。王恢等兵三万。闻单于不与汉合。度往击辎重。必与单于精兵战。汉兵势必败。则以便宜罢兵。皆无功。\\
\hspace*{2em}天子怒王恢不出击单于辎重。擅引兵罢也。恢曰。﹁始约虏入马邑城。兵与单于接。而臣击其辎重。可得利。今单于闻。不至而还。臣以三万人众不敌。礻是取辱耳。臣固知还而斩。然得完陛下士三万人。﹂于是下恢廷尉。廷尉当恢逗桡。当斩。恢私行千金丞相蚡。蚡不敢言上。而言于太后曰。﹁王恢首造马邑事。今不成而诛恢。是为匈奴报仇也。﹂上朝太后。太后以丞相言告上。上曰。﹁首为马邑事者。恢也。故发天下兵数十万。从其言。为此。且纵单于不可得。恢所部击其辎重。犹颇可得。以慰士大夫心。今不诛恢。无以谢天下。﹂于是恢闻之。乃自杀。\\
\hspace*{2em}安国为人多大略。智足以当世取合。而出于忠厚焉。贪嗜于财。所推举皆廉士。贤于己者也。于梁举壶遂。臧固。郅他。皆天下名士。士亦以此称慕之。唯天子以为国器。安国为御史大夫四岁馀。丞相田蚡死。安国行丞相事。奉引堕车蹇。天子议置相。欲用安国。使使视之。蹇甚。乃更以平棘侯薛泽为丞相。安国病免数月。蹇愈。上复以安国为中尉。岁馀。徙为卫尉。\\
\hspace*{2em}车骑将军卫青击匈奴。出上谷。破胡茏城。将军李广为匈奴所得。复失之。公孙敖大亡卒。皆当斩。赎为庶人。明年。匈奴大入边。杀辽西太守。及入雁门。所杀略数千人。车骑将军卫青击之。出雁门。卫尉安国为材官将军。屯于渔阳。安国捕生虏。言匈奴远去。即上书言方田作时。请且罢军屯。罢军屯月馀。匈奴大入上谷。渔阳。安国壁乃有七百馀人。出与战。不胜。复入壁。匈奴虏略千馀人及畜产而去。天子闻之。怒。使使责让安国。徒安国益东。屯右北平。是时匈奴虏言当入东方。\\
\hspace*{2em}安国始为御史大夫及护军。后稍斥疏。下迁。而新幸壮将军卫青等有功。益贵。安国既疏远。默默也。将屯又为匈奴所欺。失亡多。甚自愧。幸得罢归。乃益东徙屯。意忽忽不乐。数月。病欧血死。安国以元朔二年中卒。\\
\hspace*{2em}太史公曰。余与壶遂定律历。观韩长孺之义。壶遂之深中隐厚。世之言梁多长者。不虚哉。壶遂官至詹事。天子方倚以为汉相。会遂卒。不然。壶遂之内廉行脩。斯鞠躬君子也。\\

% ===== 索隐述赞（司马贞《史记索隐》）· 统一样式 =====
% 题头「索隐述赞曰」五字并入反白框；正文每四句一列，定宽盒 \hfil 居列中；
% \fontsize{24}{41} 列距对齐全书界行（121+41k）；\clearpage 使述赞独占一页。
\clearpage
\shuzanhead{索隐述赞曰}\linebreak
{\shuzanbody\fontsize{24}{41}\selectfont
\hbox to \hsize{\hfil 安国忠厚。初为梁将。因事坐法。免徒起相。\hfil}\linebreak
\hbox to \hsize{\hfil 死灰更然。生虏失防。推贤见重。贿金贻谤。\hfil}\linebreak
\hbox to \hsize{\hfil 雪泣悟主。臣节可亮。\hfil}%
\par}
\end{document}
